(1) 微小弦元横向张力合力为 $T\partial_x^2y\,dx$，黏滞力为 $-\gamma\partial_ty\,dx$，牛顿第二定律给 $\mu y_{tt}+\gamma y_t=Ty_{xx}$。两端固定为 $y(0,t)=y(L,t)=0$。$\mu$ 的单位是质量每长度，而 $\gamma$ 的单位是质量每长度每时间，两者不可混用。
(2) 用正弦本征函数代入方程得到 $\ddot q_n+(\gamma/\mu)\dot q_n+\omega_n^2q_n=0$，其中 $\omega_n=(n\pi/L)\sqrt{T/\mu}$。定义 $\alpha=\gamma/(2\mu)$，欠阻尼条件是 $\omega_n>\alpha$，此时振荡角频率 $\omega_{dn}=\sqrt{\omega_n^2-\alpha^2}$，不是原来无阻尼的 $\omega_n$。
(3) 该模态方程的通解为 $q_n=e^{-\alpha t}[C\cos(\omega_{dn}t)+D\sin(\omega_{dn}t)]$。施加 $q_n(0)=A$ 得 $C=A$，施加 $\dot q_n(0)=0$ 得 $D=\alpha A/\omega_{dn}$。故 $q_n(t)=Ae^{-\alpha t}[\cos(\omega_{dn}t)+(\alpha/\omega_{dn})\sin(\omega_{dn}t)]$，同时满足初始位移与速度约束。
(4) 把波方程乘 $y_t$ 并沿 $0<x<L$ 积分，利用固定端边界可得 $dE/dt=-\gamma\int_0^Ly_t^2dx\le0$，其中 $E=\int_0^L[\mu y_t^2+Ty_x^2]dx/2$。弱阻尼下模态振幅包络为 $e^{-\alpha t}$，能量包络近似 $e^{-2\alpha t}=e^{-\gamma t/\mu}$，寿命 $\tau_E=\mu/\gamma$，品质因子 $Q_n\simeq\omega_n/(2\alpha)=\mu\omega_n/\gamma$。
